\section{One original centre family for free jets and tangent symbols}

\begin{lemma}[Actual zero-translation generators]
\label{hir:actual-coupled-zero-translation-generators}
\lean{Hironaka.HironakaActualCentreCoupledSymbols.translatedIdeal_zero_eq_span}
\leanok
For a family of holomorphic functions near the origin, its genuine
translated germ ideal at zero equals the entire ideal spanned by the actual
untranslated origin germs. The equality is independent of the proofs of
holomorphicity used to form those germs.
\end{lemma}
\begin{proof}
The translated representative of each function at zero is the same function:
addition of zero changes no argument. Thus the two representatives agree on
a neighbourhood of zero, hence give identical holomorphic germs. Equality of
the complete generating ranges gives equality of their whole ideal spans.
\end{proof}

\begin{theorem}[Coupled all-order freeness and initial tangent stability]
\label{hir:actual-coupled-centre-free-jets-initial-symbols}
\lean{Hironaka.HironakaActualCentreCoupledSymbols.exists_original_hs_constant_centre_coupled_symbols}
\leanok
\uses{hir:actual-coupled-zero-translation-generators,hir:actual-hs-constant-centre-all-order-principal-parts-free,hir:actual-original-centre-frame-first-symbol-tangent-closure,hir:actual-original-first-symbol-whole-initial-pderiv-transfer}
Let $C$ be regular of dimension $d$ at $x$ in a complex manifold of
dimension $d+t$, and let $A$ be a coherent ideal with its whole intrinsic
Hilbert--Samuel sequence constant along $C$. These original data supply
one flattening chart, one neighbourhood and one actual finite generator
family $H$, with origin ideal $I_0$ and actual joint ideal $J$.
They simultaneously prove:
\[
 I_0=(H_1,\ldots,H_r),\qquad
 R/(J+N^q)\text{ finite projective and free over }B\quad(q>0),
\]
and stability of the whole genuine Taylor initial ideal of $I_0$ under
every tangent coordinate derivative. Both formal-series stability and its
literal polynomial version hold. Here $B$ is the actual tangent germ ring,
$R$ the unchanged joint germ ring and $N$ its entire normal-coordinate ideal.
\end{theorem}
\begin{proof}
Apply the original all-order principal-parts producer once. It supplies one
chart and generator family, all original stalk identifications, and at each
order an actual invertible pivot minor with its constant rank, together
with finite projectivity and freeness of the whole quotient. Retain all
these witnesses, rather than choosing a second flattening for tangent
stability. The zero-translation lemma identifies the same origin ideal
with the span of the same origin germs.

At order $q>0$, the same pivot and rank data give the projected tangent
derivative identity. For an actual $g\in I_0\cap\mathfrak m^q$ it says
$\partial_i g\in I_0+\mathfrak m^q$. The genuine Taylor derivative
identity and homogeneous component argument therefore show that the
entire initial ideal is stable under the same tangent derivative. The
zero-base Taylor coordinate identification retains exactly the original
unshifted germ coordinates. Finally the polynomial-to-series derivative
identity restricts this whole-series conclusion to every polynomial.
All degrees use the original family fixed before the choice of $q$.
\end{proof}
